\documentclass[12pt, a4paper]{article}
\usepackage[UTF8]{ctex}

\usepackage{amsmath, amssymb, amsthm}
\usepackage{geometry}
\usepackage{fancyhdr}
\usepackage{enumitem}
\usepackage{xcolor}
\usepackage{hyperref}

\geometry{margin=1in}
\pagestyle{fancy}
\fancyhf{}
\rhead{Mathematics}
\lhead{Li Shengdan}
\rfoot{第 \thepage 页}

\newcommand{\problem}[1]{\subsection*{题 #1}}
\newcommand{\solution}{\paragraph{解.}}

\definecolor{solutioncolor}{RGB}{0, 100, 0}

\begin{document}

\begin{center}
    {\LARGE\bfseries Homework Solutions} \\[0.5em]
    {\large Mathematics} \\[0.3em]
    Li Shengdan \quad | \quad 2026-10-05
\end{center}

\hrule
\vspace{1em}
\problem{PartA1}

Find the derivative of $x^3 + 2x^2 - 5x + 1$.

\solution

对多项式逐项求导：$(x^3)' = 3x^2$

$(2x^2)' = 4x$

$(-5x)' = -5$

$(1)' = 0$

合并得 $3x^2 + 4x - 5$

\textbf{Answer:}
\[\boxed{3x^2 + 4x - 5}\]

\vspace{1em}
\hrule
\vspace{1em}

\problem{PartA2}

Evaluate the integral $\int x^2 \sin(x) dx$.

\solution

使用分部积分，令 $u = x^2$，$dv = \sin(x)dx$，则 $du = 2x\,dx$，$v = -\cos(x)$

$\int x^2\sin(x)dx = -x^2\cos(x) + 2\int x\cos(x)dx$

对 $\int x\cos(x)dx$ 再次分部积分：$u = x$，$dv = \cos(x)dx$，得 $x\sin(x) - \int \sin(x)dx = x\sin(x) + \cos(x)$

代回得 $-x^2\cos(x) + 2x\sin(x) + 2\cos(x) + C$

\textbf{Answer:}
\[\boxed{-x^2\cos(x) + 2x\sin(x) + 2\cos(x) + C}\]

\vspace{1em}
\hrule
\vspace{1em}

\problem{PartA3}

Compute the limit $\lim_{x \to 0} \frac{\sin(x)}{x}$.

\solution

$\lim_{x \to 0} \frac{\sin{\left(x \right)}}{x} = 1$

\textbf{Answer:}
\[\boxed{1}\]

\vspace{1em}
\hrule
\vspace{1em}

\problem{PartA4}

Find the derivative of $e^{x^2} \cdot \ln(x)$.

\solution

使用乘积法则：$(uv)' = u'v + uv'$，其中 $u = e^{x^2}$，$v = \ln(x)$

$u' = e^{x^2} \cdot 2x$（链式法则）

$v' = \frac{1}{x}$

因此导数为 $2x e^{x^2}\ln(x) + \frac{e^{x^2}}{x}$

\textbf{Answer:}
\[\boxed{2x e^{x^2}\ln(x) + \frac{e^{x^2}}{x}}\]

\vspace{1em}
\hrule
\vspace{1em}

\problem{PartA5}

Evaluate $\int_0^1 \frac{1}{1+x^2} dx$.

\solution

$\frac{1}{1+x^2}$ 的原函数为 $\arctan(x)$

由牛顿-莱布尼茨公式：$\int_0^1 \frac{1}{1+x^2}dx = \arctan(1) - \arctan(0)$

$\arctan(1) = \frac{\pi}{4}$，$\arctan(0) = 0$

故积分值为 $\frac{\pi}{4}$

\textbf{Answer:}
\[\boxed{\frac{\pi}{4}}\]

\vspace{1em}
\hrule
\vspace{1em}

\problem{PartB6}

Solve the equation $x^2 - 5x + 6 = 0$.

\solution

因式分解：$x^2 - 5x + 6 = (x-2)(x-3)$

令 $(x-2)(x-3) = 0$

得 $x = 2$ 或 $x = 3$

\textbf{Answer:}
\[\boxed{x = 2 \text{ 或 } x = 3}\]

\vspace{1em}
\hrule
\vspace{1em}

\problem{PartB7}

Solve $x^3 - 6x^2 + 11x - 6 = 0$.

\solution

试根 $x=1$：$1 - 6 + 11 - 6 = 0$，故 $(x-1)$ 是因子

多项式除法：$x^3 - 6x^2 + 11x - 6 = (x-1)(x^2 - 5x + 6)$

继续分解：$x^2 - 5x + 6 = (x-2)(x-3)$

因此 $x^3 - 6x^2 + 11x - 6 = (x-1)(x-2)(x-3) = 0$

得 $x = 1$，$x = 2$ 或 $x = 3$

\textbf{Answer:}
\[\boxed{x = 1, x = 2 \text{ 或 } x = 3}\]

\vspace{1em}
\hrule
\vspace{1em}

\problem{PartC8}

Compute the determinant of the matrix $A = \begin{pmatrix} 1 & 2 \\ 3 & 4 \end{pmatrix}$.

\solution

矩阵 $A = \left[\begin{matrix}1 & 2\\3 & 4\end{matrix}\right]$

按行列式定义展开（SymPy 符号计算，逐步约化）

$\det(A) = -2$

\textbf{Answer:}
\[\boxed{-2}\]

\vspace{1em}
\hrule
\vspace{1em}

\problem{PartC9}

Find the eigenvalues of $A = \begin{pmatrix} 2 & 1 \\ 1 & 2 \end{pmatrix}$.

\solution

矩阵 $A = \left[\begin{matrix}2 & 1\\1 & 2\end{matrix}\right]$

特征方程 $\det(A - \lambda I) = 0$，展开求解：

$\lambda = 3$, $\lambda = 1$

特征值: $\lambda_{1} = 3, \; \lambda_{2} = 1$

\textbf{Answer:}
\[\boxed{\lambda_{1} = 3, \; \lambda_{2} = 1}\]

\vspace{1em}
\hrule
\vspace{1em}

\problem{PartD10}

Solve the ODE $y' + 2y = e^{-x}$.

\solution

解析得到微分方程: $2 y{\left(x \right)} + \frac{d}{d x} y{\left(x \right)} - e^{- x} = 0$

SymPy dsolve 求解（1 阶）

通解: $y = \left(C_{1} e^{- x} + 1\right) e^{- x}$（含任意常数 C1）

\textbf{Answer:}
\[\boxed{y = \left(C_{1} e^{- x} + 1\right) e^{- x}}\]

\vspace{1em}
\hrule
\vspace{1em}


\vspace{2em}
\noindent\textit{Auto-generated: 10/10 problems solved by SymPy.}
\end{document}
